% Methods documentation for the anisotropy package.
% Compile: pdflatex README.tex  (twice if references needed)
\documentclass[11pt,a4paper]{article}

\usepackage[margin=1in]{geometry}
\usepackage{amsmath,amssymb,bm}
\usepackage{booktabs}
\usepackage{hyperref}
\usepackage{listings}
\usepackage{microtype}

\hypersetup{colorlinks=true, linkcolor=blue!60!black, citecolor=blue!60!black}

\title{Protein SAS Mesh, Shape Descriptors, and Patch Parameterization\\
\large \texttt{anisotropy} package (implementation reference)}
\author{Cryo-EM toys / thesis workflow}
\date{\today}

\begin{document}
\maketitle

\begin{abstract}
This note defines the geometric and chemical quantities computed by the
\texttt{anisotropy} Python package: construction of a solvent-accessible surface (SAS)
triangle mesh from a PDB structure, global shape descriptors (radius of gyration and
asphericity) from that mesh, and patch-resolved feature vectors for interfacial
modeling. All formulas match the current source files
\texttt{sasa.py}, \texttt{mesh.py}, \texttt{shape.py}, and \texttt{patches.py}.
\end{abstract}

\tableofcontents
\newpage

%=============================================================================
\section{Input and coordinate frame}
%=============================================================================

A structure is read from a PDB file as a list of heavy atoms (ATOM records; HETATM
optional). Each atom $i$ has position $\mathbf{c}_i\in\mathbb{R}^3$ ({\AA}) and
van der Waals radius $r_i^{\mathrm{vdW}}$ (Bondi radii). A water probe of radius
$r_{\mathrm{p}}$ (default $r_{\mathrm{p}}=1.4$\,{\AA}) defines the solvent-accessible
radius $R_i = r_i^{\mathrm{vdW}} + r_{\mathrm{p}}$.

All meshes and descriptors use the same laboratory Cartesian frame as the PDB.

%=============================================================================
\section{Solvent-accessible surface and mesh formation}
%=============================================================================

\subsection{Algebraic SAS distance field (exact backend)}

For any grid point $\mathbf{x}$, the \emph{algebraic} SAS distance function is
\begin{equation}
  \Phi(\mathbf{x})
  = \min_i \bigl( \|\mathbf{x}-\mathbf{c}_i\| - R_i \bigr).
  \label{eq:phi-exact}
\end{equation}
Negative values lie inside the union of solvent-accessible spheres; the SAS
\emph{surface} is the zero level set $\{\mathbf{x}:\Phi(\mathbf{x})=0\}$.

The exact backend samples $\Phi$ on a uniform grid of spacing $h$ covering the
structure bounding box (with padding), by evaluating \eqref{eq:phi-exact} at every
voxel center. If the grid would exceed a maximum edge length, $h$ is increased
automatically.

\subsection{Voxel SAS signed distance field (fast backend)}

For large structures ($N\gtrsim 1.2\times 10^4$ atoms in the \texttt{auto} mode, or
when \texttt{method=voxel} is selected), an equivalent surface is built by:
\begin{enumerate}
  \item Rasterizing the union of spheres of radius $R_i$ onto a boolean grid
        (voxels marked \emph{inside} if covered by any sphere).
  \item Computing Euclidean distance transforms:
        \begin{align}
          d_{\mathrm{out}}(\mathbf{x}) &= \text{EDT}(\text{outside}), \\
          d_{\mathrm{in}}(\mathbf{x})  &= \text{EDT}(\text{inside}),
        \end{align}
        and forming the signed distance
        \begin{equation}
          \mathrm{SDF}(\mathbf{x}) =
          \begin{cases}
            -\,d_{\mathrm{in}}(\mathbf{x}), & \mathbf{x}\ \text{inside},\\[2pt]
            \ \ d_{\mathrm{out}}(\mathbf{x}), & \mathbf{x}\ \text{outside}.
          \end{cases}
        \end{equation}
  \item Optionally applying a Gaussian smooth to the field ($\sigma\approx 0.75$
        voxels) for numerical stability.
\end{enumerate}
The mesh surface is extracted at $\mathrm{SDF}=0$, which approximates the SAS
envelope used in \eqref{eq:phi-exact}.

\subsection{Isosurface extraction (marching cubes)}

Given a scalar field $F$ on a grid ($F=\Phi$ or $F=\mathrm{SDF}$), the
\emph{marching-cubes} algorithm (via \texttt{skimage.measure.marching\_cubes})
extracts a triangle mesh $\mathcal{M}=(\mathcal{V},\mathcal{F})$:
\begin{itemize}
  \item vertices $\mathbf{v}_m\in\mathcal{V}$, $m=1,\ldots,V$;
  \item triangular faces $f_\ell=(i_1,i_2,i_3)\in\mathcal{F}$, $\ell=1,\ldots,F$.
\end{itemize}
Vertex positions are shifted from grid indices to {\AA} using the grid origin and
spacing $h$.

\subsection{Iterative refinement}

The routine \texttt{fit\_iterative\_mesh} repeats surface extraction at a sequence
of grid spacings (defaults $h\in\{2.5,\,1.5,\,1.0\}$\,\AA):
\begin{enumerate}
  \item build $F$ at resolution $h$;
  \item extract $\mathcal{M}$ at level $F=0$;
  \item apply Laplacian smoothing on vertex positions (uniform weights on the mesh
        graph), with iteration count increasing on finer passes (defaults
        $2,3,4,\ldots$).
\end{enumerate}
Each pass replaces the previous mesh; the finest pass is retained. Smoothing uses
\begin{equation}
  \mathbf{v}_m^{(t+1)}
  = \mathbf{v}_m^{(t)}
  + \lambda \left(
      \frac{1}{|\mathcal{N}(m)|}\sum_{n\in\mathcal{N}(m)} \mathbf{v}_n^{(t)}
      - \mathbf{v}_m^{(t)}
    \right),
  \label{eq:laplacian}
\end{equation}
with $\lambda=0.45$ and $\mathcal{N}(m)$ the set of vertices adjacent to $m$ on the
mesh.

\subsection{Triangle geometry}

For face $\ell$ with vertices $\mathbf{v}_1,\mathbf{v}_2,\mathbf{v}_3$, the
implementation uses
\begin{align}
  \mathbf{A}_\ell &= \tfrac{1}{2}(\mathbf{v}_2-\mathbf{v}_1)\times(\mathbf{v}_3-\mathbf{v}_1),\\
  a_\ell &= \|\mathbf{A}_\ell\|, \qquad
  \hat{\mathbf{n}}_\ell = \frac{\mathbf{A}_\ell}{\|\mathbf{A}_\ell\|},\\
  \mathbf{r}_\ell &= \tfrac{1}{3}(\mathbf{v}_1+\mathbf{v}_2+\mathbf{v}_3),
\end{align}
where $a_\ell$ is the face area and $\hat{\mathbf{n}}_\ell$ the outward unit normal
(orientation follows the vertex ordering from marching cubes).

%=============================================================================
\section{Global shape descriptors from the mesh}
%=============================================================================

Shape descriptors are computed from the \textbf{mesh vertex sample} of the SAS
surface, not from atomic centers alone. This matches the idea that global
anisotropy is a property of the \emph{surface envelope}.

\subsection{Vertex weights}

Let $\mathbf{v}_m$ be mesh vertices and $\mathcal{F}$ the set of triangular faces.
Each vertex receives a weight $w_m$ proportional to one third of the sum of areas of
incident triangles (then floored and normalized):
\begin{equation}
  \tilde{w}_m = \frac{w_m}{\sum_j w_j}, \qquad \sum_m \tilde{w}_m = 1.
\end{equation}

\subsection{Center and gyration tensor}

The reference center is
\begin{equation}
  \mathbf{r}_c = \sum_{m=1}^{V} \tilde{w}_m\, \mathbf{v}_m.
\end{equation}
Define centered coordinates $\boldsymbol{\rho}_m = \mathbf{v}_m - \mathbf{r}_c$.
The symmetric \textbf{gyration tensor} (second moment matrix) is
\begin{equation}
  \mathbf{Q}
  = \sum_{m=1}^{V} \tilde{w}_m\, \boldsymbol{\rho}_m \boldsymbol{\rho}_m^{\mathsf{T}},
  \qquad Q_{ab} = \sum_m \tilde{w}_m\, \rho_{m,a}\,\rho_{m,b}.
  \label{eq:gyration}
\end{equation}
Units are {\AA}$^2$ because $\tilde{w}_m$ are dimensionless.

Diagonalize $\mathbf{Q}$:
\begin{equation}
  \mathbf{Q}\,\mathbf{e}_n = \lambda_n\,\mathbf{e}_n,
  \qquad n\in\{1,2,3\},
\end{equation}
with eigenvalues ordered $\lambda_1\ge\lambda_2\ge\lambda_3\ge 0$.
The principal axes are unit vectors $\mathbf{e}_1,\mathbf{e}_2,\mathbf{e}_3$ (rows of
the eigenvector matrix in code).

\subsection{Radius of gyration}

The implementation defines
\begin{equation}
  \boxed{
    R_g = \sqrt{\lambda_1+\lambda_2+\lambda_3}
        = \sqrt{\mathrm{tr}(\mathbf{Q})}.
  }
  \label{eq:Rg}
\end{equation}
Thus $R_g^2$ is the total weighted mean-square displacement of surface vertices about
$\mathbf{r}_c$ (summed over all principal directions). This is the gyration radius
associated with $\mathbf{Q}$, not the radius of a sphere with equal volume.

Effective semi-axis lengths reported in code:
$L_n = 2\sqrt{\lambda_n}$.

\subsection{Asphericity (not Wadell sphericity)}

The code computes \textbf{asphericity} $\mathcal{A}$, \emph{not} the classical
Wadell sphericity $\psi = \pi^{1/3}(6V)^{2/3}/A$.
Let $\Sigma=\lambda_1+\lambda_2+\lambda_3$. Then
\begin{equation}
  \boxed{
    \mathcal{A}
    = 1 - \frac{3\lambda_3}{\Sigma}
    = 1 - \frac{3\lambda_3}{\lambda_1+\lambda_2+\lambda_3}.
  }
  \label{eq:asphericity}
\end{equation}
\begin{itemize}
  \item $\mathcal{A}=0$ when $\lambda_1=\lambda_2=\lambda_3$ (isotropic second moments).
  \item $\mathcal{A}\to 1$ when $\lambda_3\ll\lambda_1,\lambda_2$ (strong elongation
        in one principal direction).
\end{itemize}
If you need \emph{sphericity} in the sense $\psi\in(0,1]$ with $\psi=1$ for a perfect
sphere, use a separate definition based on mesh volume and area; it is \textbf{not}
currently implemented.

\subsection{Prolateness (also reported)}

The package additionally reports
\begin{equation}
  \mathcal{P}
  = \frac{3\lambda_1-\Sigma}{2\Sigma}
    - \frac{3\lambda_3-\Sigma}{2\Sigma},
\end{equation}
as a signed measure related to prolate versus oblate ellipsoid character.

\subsection{Curvature on patches (used later)}

Per-vertex mean curvature $H$ and Gaussian curvature $K$ are obtained from the mesh
via VTK/PyVista (\texttt{curvature} filters) and averaged over patch vertices.

%=============================================================================
\section{Patch segmentation and parameterization}
%=============================================================================

\subsection{Patch decomposition}

Faces of $\mathcal{M}$ are grouped into patches by \textbf{region growing} on the
dual graph: two adjacent faces belong to the same patch if the angle between their
unit normals satisfies
\begin{equation}
  \hat{\mathbf{n}}_\ell \cdot \hat{\mathbf{n}}_{\ell'} \ge \cos\theta_{\max},
\end{equation}
with default $\theta_{\max}=25^\circ$. Seeds are processed in descending face area.
Patches with total area below $a_{\min}$ (default $50$\,\AA$^2$) are merged into a
neighboring patch.

Each patch $f$ has a set of face indices $\mathcal{F}_f$ and inherits triangle
geometry from Section~2.5.

\subsection{Patch feature vector}

The modeling proposal assigns each patch a vector
\begin{equation}
  \mathbf{x}_f
  = \bigl(
    a_f,\,
    \mathbf{r}_f,\,
    \hat{\mathbf{n}}_f,\,
    H_f,\,
    K_f,\,
    q_f,\,
    \phi_f,\,
    \mathrm{pK}_{a,f},\,
    h_f,\,
    p_f,\,
    b_f,\,
    \boldsymbol{\mu}_f,\,
    u_f
  \bigr).
\end{equation}
Table~\ref{tab:patch} lists definitions implemented in \texttt{patches.py}.

\begin{table}[ht]
  \centering
  \caption{Patch descriptors $\mathbf{x}_f$ (implementation).}
  \label{tab:patch}
  \begin{tabular}{@{}llp{7.2cm}@{}}
    \toprule
    Symbol & Code name & Definition \\
    \midrule
    $a_f$ & \texttt{area}
      & $a_f=\sum_{\ell\in\mathcal{F}_f} a_\ell$ \\
    $\mathbf{r}_f$ & \texttt{centroid}
      & Area-weighted mean of face centroids $\mathbf{r}_\ell$ \\
    $\hat{\mathbf{n}}_f$ & \texttt{normal}
      & Area-weighted mean of $\hat{\mathbf{n}}_\ell$, then normalized \\
    $H_f$, $K_f$ & \texttt{mean\_curvature}, \texttt{gaussian\_curvature}
      & Mean over unique vertices of faces in patch $f$ \\
    $q_f$ & \texttt{charge}
      & $\sum_{i\in\mathcal{I}_f} q_i(pH)$; see below \\
    $\phi_f$ & \texttt{potential}
      & $\sum_{i\in\mathcal{I}_f} k\, q_i / \max(\|\mathbf{c}_i-\mathbf{r}_f\|,1)$ \\
    $\mathrm{pK}_{a,f}$ & \texttt{pka\_acid}
      & Mean residue p$K_\mathrm{a}$ over titratable types in patch \\
    $h_f$ & \texttt{hydropathy}
      & $\bar{h}_{\mathrm{KD}}\,\|\hat{\mathbf{n}}_f\|$ (directional hydrophobicity) \\
    $p_f$ & \texttt{polar\_density}
      & (\# polar residues in patch)/(\# atoms in patch) \\
    $b_f$ & \texttt{hbond\_score}
      & (donors$+$acceptors per residue)/(\# atoms in patch) \\
    $\boldsymbol{\mu}_f$ & \texttt{dipole}
      & $\sum_{i\in\mathcal{I}_f} q_i\,(\mathbf{c}_i-\mathbf{r}_f)$ \\
    $u_f$ & \texttt{softness}
      & Mean B-factor$/100$ if PDB B-factors exist; else curvature proxy \\
    \bottomrule
  \end{tabular}
\end{table}

\subsection{Atom-to-patch assignment}

Let $\mathbf{r}_\ell$ be face centroids. Each atom $i$ is assigned to the patch of
its nearest face centroid:
\begin{equation}
  f(i) = \arg\min_\ell \|\mathbf{c}_i - \mathbf{r}_\ell\|,
  \qquad
  \mathcal{I}_f = \{ i : f(i)=f \}.
\end{equation}
All chemical sums run over $i\in\mathcal{I}_f$.

\subsection{Charge and protonation at pH}

For residue type $R$, a simplified Henderson--Hasselbalch model gives fractional
charge $q_R(pH)$ (acidic residues negative, basic positive). Each atom carries a
scaled partial charge proxy $q_i(pH)$ derived from its residue. The patch charge is
$q_f=\sum_{i\in\mathcal{I}_f} q_i(pH)$.

\subsection{Electrostatic potential proxy}

With Coulomb prefactor $k=332.0636$\,(kcal\,mol$^{-1}$\,\AA\,e$^{-2}$) in code,
\begin{equation}
  \phi_f = \sum_{i\in\mathcal{I}_f} \frac{k\, q_i(pH)}{\max(\|\mathbf{c}_i-\mathbf{r}_f\|,1\,\mathrm{\AA})}.
\end{equation}
This is a \emph{coarse} pairwise sum, not a Poisson--Boltzmann solve (APBS).

\subsection{Hydropathy and interface direction}

Kyte--Doolittle hydropathy $h^{\mathrm{KD}}_R$ is assigned per residue ($h>0$
hydrophobic). For patch $f$,
\begin{equation}
  \bar{h}_{\mathrm{KD}} = \frac{1}{|\mathcal{I}_f|}\sum_{i\in\mathcal{I}_f} h^{\mathrm{KD}}_{R(i)},
  \qquad
  h_f = \bar{h}_{\mathrm{KD}}\,\|\hat{\mathbf{n}}_f\|.
\end{equation}
Thus hydrophobic character is weighted by how much the patch normal participates
in the outward direction (directional presentation to the interface).

\subsection{Hydrogen-bond channel}

Per-residue donor and acceptor counts $(d_R,a_R)$ are tabulated. Then
\begin{equation}
  b_f = \frac{\sum_{i\in\mathcal{I}_f}(d_{R(i)}+a_{R(i)})}{|\mathcal{I}_f|}.
\end{equation}

\subsection{Softness}

If crystallographic B-factors $B_i$ are present in the PDB,
\begin{equation}
  u_f = \frac{1}{100\,|\mathcal{I}_f|}\sum_{i\in\mathcal{I}_f} B_i.
\end{equation}
Otherwise a heuristic combines mean curvature and polarity:
\begin{equation}
  u_f = \min\!\Bigl(1,\; 2|H_f| + 0.3(1-\min(p_f,1))\Bigr).
\end{equation}

%=============================================================================
\section{Software entry points}
%=============================================================================

\begin{itemize}
  \item \texttt{fit\_protein\_mesh.py} --- PDB $\to$ SAS mesh (PLY); global
        $R_g$, $\mathcal{A}$.
  \item \texttt{parameterize\_mesh.py} --- mesh $+$ PDB $\to$
        \texttt{patch\_features.npz} and optional JSON.
  \item Python API: \texttt{fit\_iterative\_mesh}, \texttt{shape\_anisotropy\_from\_mesh},
        \texttt{parameterize\_mesh}.
\end{itemize}

\paragraph{Recommended interpreter.}
Use the toys virtual environment:
\texttt{toys/.venv/Scripts/python.exe} (Python~3.12).
Plain \texttt{python} on some Windows installs resolves to MGLTools Python~2.7.

%=============================================================================
\section{Summary of key formulas}
%=============================================================================

\begin{align}
  \Phi(\mathbf{x}) &= \min_i(\|\mathbf{x}-\mathbf{c}_i\|-R_i),
  &\text{(exact SAS)} \\[4pt]
  \mathbf{Q} &= \sum_m \tilde{w}_m\,\boldsymbol{\rho}_m\boldsymbol{\rho}_m^{\mathsf{T}},
  &\text{(gyration tensor)} \\[4pt]
  R_g &= \sqrt{\mathrm{tr}(\mathbf{Q})},
  &\text{(radius of gyration)} \\[4pt]
  \mathcal{A} &= 1 - \frac{3\lambda_3}{\lambda_1+\lambda_2+\lambda_3},
  &\text{(asphericity)}.
\end{align}

\end{document}
